\documentclass[11pt]{article}
\usepackage{amsmath,amssymb}
\title{Cohort-by-date fixed effects from observed outcome differences}
\author{}
\date{}
\begin{document}
\maketitle

Let \(O_i\) denote the ordered dates with an observed outcome for study unit \(i\).
Write \(m_i=|O_i|\). For \(m_i\ge 1\), collect the observed outcomes in \(y_i\)
and write their linear model as
\[
  y_i = \mathbf{1}_{m_i}\alpha_i + Z_i\beta + u_i.
\]
Rows of \(Z_i\) pick a coefficient slot. There is one common slot \(\lambda_t\)
per date. A saturated cohort-date cell instead receives its own slot
\(\gamma_{g,t}\). The treatment contrast for a post-treatment cell is
\(\tau_{g,t}=\gamma_{g,t}-\lambda_t\). With never-treated controls the
cell \((g,g-1)\) uses the common slot and other pre-treatment cells have
separate deviations. With not-yet-treated controls all untreated cells use
common slots, imposing common pre-treatment changes across cohorts.

\paragraph{Eliminating the unit intercept.}
For \(m_i\ge 2\) let \(B_i\) form successive differences along the observed
dates, allowing calendar gaps. Its \(m_i-1\) rows are linearly independent,
its kernel is the span of \(\mathbf{1}_{m_i}\), and its row space is the
orthogonal complement of that vector. Therefore
\[
  B_i' (B_i B_i')^{-1} B_i
  = M_i
  = I_{m_i}-\frac{\mathbf{1}_{m_i}\mathbf{1}_{m_i}'}{m_i}.
\]
Both sides are the orthogonal projection onto the same row space. The
within-unit least-squares objective is thus exactly
\[
  \sum_{i:m_i\ge 2}
  (B_i y_i-B_i Z_i\beta)'(B_i B_i')^{-1}(B_i y_i-B_i Z_i\beta).
\]
It uses \(m_i-1\) differences regardless of how many additional pairwise
differences can be formed. It weights correlated differences jointly.
A one-date unit contributes zero after elimination. Units with no observed
outcomes have no rows but can remain in the unit frame.

\paragraph{Identification.}
The Gram matrix is \(H=\sum_i Z_i'M_iZ_i\).
For any vector \(v\), \(v'Hv=\sum_i\lVert M_i Z_i v\rVert^2\).
It vanishes exactly when the slot values are constant along each unit's
observed dates. Make slots vertices and connect two slots observed by a
unit. The null space consists of component-wise constants. A contrast
\(\gamma_{g,t}-\lambda_t\) is estimable exactly when its two slots lie in
the same component. Removing one slot per component fixes the irrelevant
origins. This operation does not add identifying restrictions to contrasts.

\paragraph{Original-unit scores.}
On the gauge-fixed system, let \(C\) map coefficients to the estimable
contrasts. The contribution of unit \(i\) to the contrast estimator's
sandwich is
\[
  v_i = C H^{-1} Z_i'M_i(y_i-Z_i\widehat\beta).
\]
The retained influence function is \(N v_i\) for the full unit frame of
size \(N\). Independent-unit covariance is \(\sum_i v_iv_i'\). With
independent clusters, sum \(v_i\) within cluster before taking the outer
products. This sandwich does not use the homoskedastic model as its
variance assumption. A regression program may report an additional
degrees-of-freedom factor. Units with fewer than two outcomes have zero
score in this linear observed-outcome criterion.

The algebra equates sample computations on the same rows. Interpreting
the coefficients as population treatment effects additionally requires
the maintained parallel-trends model, treatment timing, and
representativeness of the observed changes for the relevant moment
conditions. In particular the needed mean restriction is on the
within-unit score under the observed date design. This fit does not
fit an observation probability or recover unobserved outcome changes.

\section{Time-invariant covariates}
Write \(w_i=(1,X_i')'\). Each coefficient slot now contains an intercept and
covariate slopes. The date slot holds the common calendar coefficient vector.
A free cohort-date slot holds the sum of the common vector and a cohort-date
deviation. If \(S_i\) assigns the observed rows of unit \(i\) to slots, its
regression design is \(D_i=S_i \otimes w_i'\). Unit intercepts are still
eliminated by \(M_i\), so
\[
  D_i'M_iD_i = (S_i'M_i S_i)\otimes(w_iw_i').
\]
For a common cohort and observation pattern, the first factor is fixed. The
fit can therefore accumulate the second factor over the units with that
pattern. It uses the same normal equations as the full dummy and interaction
regression.

Let \(C_{g,t}\) take the vector difference between the cohort-date and
calendar slots. The reported average effect is
\(\theta_{g,t}=\bar w_g' C_{g,t}\widehat\beta\), where \(\bar w_g\) averages
the full declared cohort frame. Support of the dates is necessary for
identification, and the covariate contrasts also need to lie in the
identified coefficient space. A finite pseudoinverse value alone does not
establish an identified contrast.

With \(N\) units and \(n_g\) in cohort \(g\), the original-unit contribution
on the package's \(N\)-scaled convention is
\[
  N\bar w_g' C_{g,t} H^{-1} D_i'M_i\widehat u_i
  +
  \frac{N 1\{G_i=g\}}{n_g}
  (w_i-\bar w_g)' C_{g,t}\widehat\beta ,
\qquad H=\sum_i D_i'M_iD_i .
\]
The inverse is evaluated on the identified gauge-fixed coordinates. The
second term comes from estimating the cohort covariate mean. A unit with
fewer than two observed outcomes can therefore contribute through its
covariates even when its regression score is zero. If a benchmark instead
holds a specified target mix fixed, only the first term is used for that
particular covariance comparison.

\section{Cohort-standardized treatment moments}
For a target cohort \(g\), write \(Q_g\) for its covariate distribution in the
study frame and \(h_{gk}(x)\) for the odds of cohort \(g\) relative to donor
\(k\) on their joint frame. For an observed clean pair \(e=(s,t)\) in cohort
\(k\), take the conditional mean untreated change to be
\(m_e^*(x)=\lambda_t(x)-\lambda_s(x)\). The same conditional function can
have a different average under each \(Q_g\).
Let \(B_{ke}=\mathbf{1}\{G=k,S_e=1\}\) indicate an observed donor change.
With working outcome prediction \(m_e(x)\), define the control moment
\[
 C_{gke}=E_{Q_g}[m_e(X)] +
   \frac{E[B_{ke} h_{gk}(X)\{\Delta_eY-m_e(X)\}]}{E[B_{ke}h_{gk}(X)]}.
\]
IPW sets \(m_e=0\). A sufficient selection-none condition is that
observed pairs preserve the conditional mean change and covariate
distribution within each relevant cohort. If \(h_{gk}\) is correct
and the cap leaves the correct odds in use, normalizing the weighted
donor distribution gives \(Q_g\). The
right-hand side then equals \(E_{Q_g}[m_e^*(X)]\). If \(m_e=m_e^*\),
the donor residual has conditional mean zero and its weighted expectation
vanishes for any admitted working odds. Together with representative
actual changes of the target cohort, these equalities give the desired
cohort-by-date contrasts.

Let \(c=E[B_{ke}h_{gk}(X)]\), \(p_g=P(G=g)\) and \(\bar r\) denote
the normalized donor residual mean. The direct original-unit term is
\[
 \frac{\mathbf{1}\{G_i=g\}}{p_g}\{m_e(X_i)-E_{Q_g}m_e(X)\}+
 \frac{B_{ke,i}h_{gk}(X_i)}{c}\{\Delta_eY_i-m_e(X_i)-\bar r\}.
\]
Each fitted nuisance contributes its influence function multiplied by
the corresponding derivative of \(C_{gke}\). For a linear outcome
nuisance this derivative is the target minus transported donor average
of its design. For the logit it is the normalized donor average of the
centered residual times the derivative of the log odds, including the
active cap rule. The graph solve applies the contrast of the weighted
least-squares loading to the original-unit moment scores. With valid
population moments, the derivative through an estimated weighting
matrix vanishes at the population solution under the usual regularity
conditions.


\section*{Reconstructing a change on the full study frame}
Let \(S_{i,s,t}\) indicate an observed change for dates \(s<t\).
Write \(q_{k,s,t}(X_i)\) for its probability in cohort \(k\),
and \(j_{k,s,t}(X_i)\) for the fitted conditional mean of the cohort's
actual change on the observed pairs. For a pair supported within that cohort,
the three reconstructed responses are
\[
 \frac{S_{i,s,t}\Delta_{i,s,t}}{q_{k,s,t}(X_i)},\qquad
 j_{k,s,t}(X_i),\qquad
 j_{k,s,t}(X_i) + \frac{S_{i,s,t}}{q_{k,s,t}(X_i)}
   \{\Delta_{i,s,t}-j_{k,s,t}(X_i)\}.
\]
Terms multiplied by \(S=0\) are zero without reading the missing outcome.
Under conditional mean representativeness, the first response has the
full-cohort conditional mean change when the probability is correct. The
second has that mean when the change model is correct. For the third response
with working functions \(q,j\) and true functions \(q_0,j_0\), the conditional
mean is \(j+(q_0/q)(j_0-j)\). It equals \(j_0\) if either function is correct
and the active clipping rule preserves the needed equality. Date-product
probabilities require the two observation indicators to be independently
distributed conditional on the selection predictors and cohort. A/B direct
pair weighting uses \(q_F(x)=\Pr(U\mid k)\rho_k(x)q_U(x)\) before the
documented clipping convention on \(q_U\).

For complete observations, the unit-eliminated normal equations are a
sum over all within-unit pairs with factor \(1/T\). The common factor
cancels. For a supported cohort-pair, substitute the reconstructed
change into the right-hand side while using the cohort's full-frame design
cross products on the left. This defines the corrected projection.
An unsupported cohort-pair does not supply a normal equation. A root offset
is fixed in each retained connected component, and any contrast with endpoints
in different components is unavailable. Treatment weighting and augmentation
instead use the reconstructed changes in the full-frame ratios of the
preceding section.

The direct unit score of the corrected projection is the design difference
times the residual of its reconstructed change. Apply the gauge-fixed
inverse and add the derivative of the right-hand side applied to the
estimating influence functions of \(q,j\). In the treatment ratios, compute
the derivative of each full-cohort mean and normalized tilt, then add its
upstream derivative loadings. A/B adds the full-frame influence of the
rate \(\pi_k=\Pr(U\mid k)\), namely
\(N/N_k\,{\bf 1}\{G_i=k\}(U_i-\pi_k)\). Its contribution to a weighted
reconstructed response is the derivative of \(1/\pi_k\). The derivative
of the probability floor is one above the floor and zero below, using
the package's stated convention at equality. Shared upstream coordinates
are summed by original unit before the sandwich or cluster aggregation.
\end{document}
